% Gerado por regressao/06_tabelas_latex.R (projeto Modelagem). Não editar à mão.
\begin{sidewaystable}
\caption{Robustez R7a — especificação principal na amostra 2013-2023}\label{tab:robustez_R7a}
\centering
\footnotesize
\setlength{\tabcolsep}{3pt}
\begin{tabular}{@{}>{\raggedright\arraybackslash}p{\dimexpr\linewidth-594pt\relax}>{\centering\arraybackslash}p{60pt}>{\centering\arraybackslash}p{60pt}>{\centering\arraybackslash}p{60pt}>{\centering\arraybackslash}p{60pt}>{\centering\arraybackslash}p{60pt}>{\centering\arraybackslash}p{60pt}>{\centering\arraybackslash}p{60pt}>{\centering\arraybackslash}p{60pt}>{\centering\arraybackslash}p{60pt}@{}}
\hline
\textbf{Variável} & \textbf{\hspace{0pt}Despesa total} & \textbf{\hspace{0pt}Saúde e Saneamento} & \textbf{\hspace{0pt}Educação e Cultura} & \textbf{\hspace{0pt}Habitação e Urbanismo} & \textbf{\hspace{0pt}Assistência Social} & \textbf{\hspace{0pt}Transporte} & \textbf{\hspace{0pt}Administração} & \textbf{\hspace{0pt}Agricultura} & \textbf{\hspace{0pt}Comunicações} \\ \hline
\multicolumn{10}{@{}l}{\textit{Modelo A (ciclo eleitoral)}} \\
Ano eleitoral & 0,0306 & 0,0446\textsuperscript{***} & 0,0030 & 0,1924\textsuperscript{***} & 0,0441 & 0,1591\textsuperscript{***} & $-$0,0106 & 0,0464 & 0,0450\textsuperscript{***} \\
 & (0,0276) & (0,0100) & (0,0477) & (0,0360) & (0,0345) & (0,0380) & (0,0156) & (0,0462) & (0,0166) \\
Ano pré-eleitoral & 0,0087 & $-$0,0112\textsuperscript{**} & 0,0297 & 0,0196 & 0,0263 & $-$0,0165 & $-$0,0005 & 0,0408 & 0,0635\textsuperscript{***} \\
 & (0,0125) & (0,0057) & (0,0198) & (0,0257) & (0,0165) & (0,0272) & (0,0076) & (0,0282) & (0,0124) \\
\hline
Observações & 59.888 & 59.747 & 59.772 & 58.890 & 59.706 & 46.153 & 59.777 & 54.492 & 10.703 \\
Municípios & 5.567 & 5.567 & 5.567 & 5.561 & 5.567 & 5.176 & 5.567 & 5.424 & 1.832 \\
R\textsuperscript{2} & 0,778 & 0,570 & 0,521 & 0,217 & 0,288 & 0,106 & 0,217 & 0,105 & 0,031 \\[6pt]
\multicolumn{10}{@{}l}{\textit{Modelo B (coalizões)}} \\
Prefeito na coalizão & 0,0033\textsuperscript{***} & $-$0,0010 & 0,0038\textsuperscript{**} & 0,0159\textsuperscript{**} & 0,0011 & 0,0165 & 0,0058\textsuperscript{**} & 0,0087 & $-$0,0071 \\
do governador & (0,0011) & (0,0020) & (0,0018) & (0,0067) & (0,0031) & (0,0104) & (0,0028) & (0,0079) & (0,0174) \\
Prefeito na coalizão & $-$0,0034\textsuperscript{**} & $-$0,0013 & $-$0,0069\textsuperscript{***} & $-$0,0181\textsuperscript{**} & $-$0,0037 & 0,0314\textsuperscript{**} & 0,0058 & 0,0006 & $-$0,0042 \\
do presidente & (0,0014) & (0,0026) & (0,0023) & (0,0087) & (0,0041) & (0,0139) & (0,0038) & (0,0104) & (0,0245) \\
\hline
Observações & 59.888 & 59.747 & 59.772 & 58.890 & 59.706 & 46.153 & 59.777 & 54.492 & 10.703 \\
Municípios & 5.567 & 5.567 & 5.567 & 5.561 & 5.567 & 5.176 & 5.567 & 5.424 & 1.832 \\
R\textsuperscript{2} & 0,330 & 0,121 & 0,146 & 0,029 & 0,060 & 0,014 & 0,061 & 0,015 & 0,009 \\
\hline
\end{tabular}
\fonte{Elaboração própria.}
\nota{Despesa e receitas per capita, em R\$ de 2025, em logaritmo (\autoref{sec:especificacao}). O coeficiente de uma dummy se lê como variação percentual aproximada (exata: 100 $\times$ (exp(b) $-$ 1)). Modelo A: efeito fixo de município, erros-padrão de Driscoll-Kraay. Modelo B: efeitos fixos de município e de ano, erros-padrão agrupados por município. Erros-padrão entre parênteses. \textsuperscript{***}p$<$0,01; \textsuperscript{**}p$<$0,05; \textsuperscript{*}p$<$0,1.}
\end{sidewaystable}
